\mode*
\section{Background: variation theory}

According to variation theory \autocite[Ch.~2]{NCOL}, each educational
objective can be divided into \emph{aspects}.
To master the objective, the learner must discern all of them; aspects not
yet discerned are \emph{critical aspects}, and misconceptions occur when a
critical aspect remains undiscerned.
Discernment requires experienced \emph{variation}: an aspect becomes visible
only when it varies against an invariant background.
Variation theory names three patterns of variation:
\emph{contrast} (the aspect varies, the rest is invariant),
\emph{generalisation} (the aspect is invariant while irrelevant features
vary), and
\emph{fusion} (several critical aspects vary together, as they do in real
use).
Phenomenography \autocite{Phenomenography,Neuman1987} supplies the
observational basis: the qualitatively different ways learners experience a
phenomenon.
In this paper, as in the companion papers
\autocite{VTProgMisconceptions,VTDebug}, the literature's documented
misconceptions play the role of those phenomenographic observations; the
instruments of \cref{sec:instruments} add the students' own accounts.
Whether the aspects so identified are the critical ones is an empirical
question, and variation theory's own answer to it is the learning study:
teach with the aspects varied, and test before and after whether the
learners now discern them \autocite[Ch.~7]{NCOL}\autocite{PangMarton2005,Kullberg2017ZDM}.
Such a test must not point out the aspects it tests---\textcquote[p.~89]{NCOL}{we
should not point out those aspects for them but let the students discern
them by themselves}---a principle the quizzes of \cref{app:quiz} follow.

\mode<presentation>{%
\begin{frame}
  \begin{block}{Variation theory \autocite{NCOL}}
    \begin{itemize}
      \item Educational objective \(=\) aspects to discern.
      \item Undiscerned critical aspect \(\rightarrow\) misconception.
      \item No discernment without variation.
      \item Patterns: contrast, generalisation, fusion.
    \end{itemize}
  \end{block}
\end{frame}%
}

A further background concept we rely on is the \emph{notional machine}: an
abstract model of the computer that novices need in order to reason about
what their actions do.
The research literature identifies the notional machine as a major challenge
in introductory programming and argues it should be an explicit learning
objective \autocite{Sorva2013NotionalMachines,DuBoulay1981BlackBox}; novice
difficulties are tied to inaccurate mental models more generally
\autocite{QianLehman2017Misconceptions,BenAri1998Constructivism}.
The shell needs a notional machine at the \emph{system} level rather than
the program-execution level: the computer runs an operating system whose
kernel is operated through a shell---the graphical desktop being one shell
and the terminal another interface to the same system.
% TODO(introtools#124, introtools#125): the du Boulay and Ben-Ari full
% texts are still to be verified before quoting specifics; see the
% provenance blocks in misconceptions.bib.
